\section{Threat Model and Security Properties}~\label{sec:sys_model}

In this section, we present our threat model and outline the privacy and security requirements our solution should satisfy.
As illustrated in Fig.~\ref{fig:overall_schema}, our system involves three entities: the client, a set of semi-trusted third parties (STTPs), and network nodes.
The client outsources header construction to STTPs, aggregates their partial headers, and injects the packet into the network.

\subsection{Threat Model}

We consider passive and active network adversaries, as well as corrupted nodes and STTPs.
Adversaries may combine these capabilities but are computationally bounded and cannot break standard cryptographic assumptions such as the hardness of the ECDLP.

\smallskip
\noindent{\textbf{Passive Network Adversary.}}
Observes packets at any node without modifying them. 
Its goal is to break unlinkability by correlating incoming and outgoing packets or exploiting statistical biases in headers.

\smallskip
\noindent{\textbf{Active Network Adversary.}}
Can modify, drop, replay, or inject packets at any point in the network.  
Its objective is to forge valid integrity tags to tamper with headers while remaining undetected. 

\smallskip
\noindent{\textbf{Corrupted Network Node.}}
Reveals its internal state, including private keys and shared secrets.
Compromising a node must not compromise the privacy or integrity of the remaining path, and unlinkability must hold as long as at least one node on the path is honest.

\smallskip
\noindent{\textbf{Corrupted STTP.}}
May attempt to use or infer information during partial header construction. 
Colluding STTPs could combine shares to reconstruct routing or shared secrets. 
Our protocol prevents this as long as at least one STTP remains honest.
We also consider cases where corrupted STTPs collude with compromised nodes, representing a more powerful combined threat.

\subsection{Security and Privacy Objectives}\label{sec:sp-objectives}

Decentralization must preserve Sphinx's cryptographic and anonymity guarantees~\cite{sphinx}, even under partial compromise and collusion.

\smallskip
\noindent\textbf{Unlinkability.}
Measuring unlinkability in anonymous communication systems is inherently challenging.
For example, Loopix~\cite{piotrowska2017loopix} proposes the \emph{expected difference in likelihood} metric, which captures sender-receiver unlinkability through network-level simulations rather than header-level resistance to cryptanalysis.
We instead focus on unlinkability at the \emph{node level}: no adversary should be able to correlate an incoming packet with its corresponding outgoing packet at any node.
This requires that header components cannot be linked without the node's shared secret.
Additionally, headers must be statistically indistinguishable from random bitstrings, a property we validate empirically using the NIST SP 800-22 test suite (Section~\ref{sec:NIST}).

\smallskip
\noindent{\textbf{Integrity-Forgery Resistance.}}
Sphinx embeds per-hop integrity tags to ensure that honest nodes are able to detect and discard altered headers.  
Any active adversary should not be able to modify a packet and forge a valid integrity tag.
In particular, no trapdoor or weakness should arise from decentralizing header construction among multiple parties.

\smallskip
\noindent{\textbf{Collusion Resistance.}} 
Additive secret sharing ensures that collusion among STTPs does not reveal routing or cryptographic secrets, as long as one STTP is honest. 
Similarly, compromised nodes should learn only local information and cannot infer the rest of the path or forge integrity, even when colluding with other nodes or STTPs.
